\section{Introduction}
\label{sec:introduction}

A common way to get more out of a large language model (LLM) at inference time is to let it check its own work. The model first produces an answer $A_0$, then critiques that answer, and finally writes a revised answer $A_1$. Self-Refine~\cite{madaan2023selfrefine} and Reflexion~\cite{shinn2023reflexion} made this loop popular, and it is now a standard building block of LLM pipelines. In this paper we study the \emph{intrinsic} version of the loop, in which the model receives no outside help: no code execution, no retrieval, and no ground-truth checker tells it whether $A_0$ was right.

Whether intrinsic self-correction actually helps is disputed. Huang et al.~\cite{huang2024cannot} found that once the oracle signal that earlier work used to decide when to stop is removed, revision on reasoning benchmarks often lowers accuracy. A survey by Kamoi et al.~\cite{kamoi2024when} concludes that reliable intrinsic self-correction has so far been shown only in narrow settings, and Stechly et al.~\cite{stechly2023gpt4} report that models are poor judges of their own solutions. The reason becomes clear once each revision is tracked individually. A revision step turns some wrong answers into right ones ($W\to R$), but it also turns some right answers into wrong ones ($R\to W$). When initial accuracy is high, there are many more correct answers that can be damaged than wrong answers that can be repaired, so even a small break rate can cancel the benefit. A single net-accuracy number hides this, because a small gain can be the sum of many fixes and many breaks.

A simple response is to revise only when the model is unsure. We call this \emph{confidence-gated self-correction}. A confidence score $S(A_0)$ is computed for the initial answer, and the revised answer is used if and only if $S(A_0)<\tau$ for a threshold $\tau$; otherwise $A_0$ is kept.

Li et al.~\cite{li2024confidence} already proposed the If-or-Else (IoE) prompt, which asks the model to keep its answer if it is confident and to revise it otherwise. In this way the decision is made inside a single prompt. Yang et al.~\cite{yang2024confidence} split self-correction into a confidence ability (keeping answers that are already correct) and a critique ability (fixing answers that are wrong). They showed that the two can trade off against each other. However, two questions are still open. First, it is not clear when gating can help at all, and how large the gain can be for a given revision step and confidence signal. Second, there is no controlled comparison of different confidence signals used as explicit threshold gates, with thresholds chosen on held-out data and with the cost of computing each signal taken into account.

This paper addresses both gaps. The study is organised around three research questions:
\begin{itemize}
    \item \textbf{RQ1 (Gating efficacy):} Does confidence gating outperform both never revising and always revising on held-out data, and how close does it come to an oracle gate?
    \item \textbf{RQ2 (Signal discrimination):} Which training-free confidence signal best separates correct from incorrect initial answers?
    \item \textbf{RQ3 (Accuracy--cost trade-off):} How does the cost of computing each signal, split into prefill and completion tokens, change the trade-off between accuracy and compute?
\end{itemize}

\paragraph{Contributions.} This paper makes three contributions.
\begin{itemize}
    \item \textbf{An analytical model of gated revision} (Section~\ref{sec:theory}). We express the accuracy of a gated system in closed form, in terms of the initial accuracy, the fix and break rates of the revision step, and the operating point of the confidence signal on its ROC curve. From this we derive when gating beats never and always revising, and show that an uninformative signal can never beat the better of the two. The model also locates the optimal threshold at a fixed slope of the ROC curve and bounds the achievable gain by the fix rate. Finally, it explains why a high AUROC alone does not guarantee a large gain.
    \item \textbf{A controlled comparison of confidence signals.} We compare three training-free signals (verbalised confidence, self-consistency agreement and P(True) self-verification) on GSM8K and HotpotQA with \modelA{}~\cite{qwen2025qwen25}. Thresholds are tuned on a development split and applied once to a test split. Differences are tested with exact McNemar tests, and the baselines include IoE, majority voting and an oracle gate.
    \item \textbf{A transition and cost analysis.} Every revision is broken down into its four possible outcomes, and cost is measured with a token model that separates prefill from completion tokens. We also inspect the $R\to W$ failures and separate changes in answer form from changes in content.
\end{itemize}
With \modelA{}, self-critique almost never fixes an error, so no gate beats never revising. Self-consistency, in contrast, identifies wrong answers well and, when used as the revision step, repairs many of them. With that revision step, gated accuracy matches the analytical model closely.

The rest of the paper is organised as follows. Section~\ref{sec:related_work} reviews related work, Section~\ref{sec:theory} presents the analytical model, Section~\ref{sec:methodology} describes the experimental setup, Section~\ref{sec:experiments} reports the results, Section~\ref{sec:analysis} analyses them in more depth. Sections~\ref{sec:discussion} and~\ref{sec:limitations} discuss implications and limitations, and Section~\ref{sec:conclusion} concludes.
